% HD-I16. GPT-6 Astra Pro, inherited effort, /root/imprimitive.
% Written and author-read 21 November 2026.
\subsection{HD-I16: exact branch-jet differential and further component ranks}
\label{sec:HD-I16}

Keep $d=3,4,6$, $e\mid d$, $n\geq5$, and the degree range of
HD-I14, but allow arbitrary local monodromies in the standard
$G(d,e,m)$ representation.  For an actual connected Nielsen component
with branch generators $\gamma_1,\ldots,\gamma_b$, let
\[
 c_j=\operatorname{codim}V^{\gamma_j},\qquad
 N=\sum_{j=1}^b(c_j-1)=2m+2n-b,
 \qquad \nu=\#\{j:c_j\geq2\}.
 \tag{I16.1}
\]
The equality follows directly from the colored-cycle identity
$c_j=\sum_{a\text{ over }j}(\ell_a-1+\boldsymbol1_{s_a\ne0})$
and the balanced signature formula of HD-I03.  A nonidentity local
monodromy has $c_j\geq1$; $c_j=1$ is exactly a reflection.
Let $g=g(Y)$ and $R$ be the reduced branch divisor of the actual
cyclic quotient $X\to Y$.

\paragraph{Statement.}
Every component with
\begin{equation}
           N+\nu\leq2m-2g+1                         \tag{I16.2}
\end{equation}
has generic period-image dimension $2n+g-1$.
In particular this holds for every component with $N\leq m-g$,
since $\nu\leq N$.  The statement includes components with
nonreflection local monodromies and does not require their
intermediate map $h:Y\to\mathbf P^1$ to be simply branched.
For arbitrary components, the explicit jet map below computes
the rank of the forgetful differential.  If $g\leq n-2$, it also
computes the period rank itself at every point.

\paragraph{1. The exact local constraints.}
Put $E=h^*T_{\mathbf P^1}$, a line bundle of degree $2m$.
At a point over the branch value $b_j$, choose local coordinates
with $h(z)=z^\ell$ and $b_j=0$.  Call the point relevant when
$\ell\geq2$ or its cyclic color $s$ is nonzero.  An unramified
uncolored point gives no branch constraint and is omitted.
Fix $(Y,R)$ and its cyclic root, and write a variation of $h$ as
\[
       h_\epsilon(z)=z^\ell+\epsilon q(z),\qquad
       q(z)=q_0+q_1z+q_2z^2+\cdots.
\]
At an uncolored ramification point the center is unmarked and
may move to $\epsilon a$.  The prescribed ramification index has
the local first-order form
\[
 \epsilon\beta+(z-\epsilon a)^\ell(1+\epsilon u(z))
   =z^\ell+\epsilon\bigl(\beta-\ell a z^{\ell-1}
                                      +z^\ell u(z)\bigr).
\]
Thus its necessary and sufficient linear constraints are
$q_1=\cdots=q_{\ell-2}=0$; the coefficient $q_{\ell-1}$ is absorbed
by the actual motion of the center.  At a colored point the point
belongs to $R$ and is fixed in this fibre of the forgetful map.
The center cannot move, so the constraints are instead
$q_1=\cdots=q_{\ell-1}=0$.  This formula includes $\ell=1$,
where there are no derivative constraints.

If there are $k_j$ relevant points over the common branch value
$b_j$, their coefficients $q_0$ must be equal, giving $k_j-1$
independent value constraints.  Their common value is the allowed
motion of $b_j$; it is not set to zero.  Consequently the number
of independent local constraints at this value is exactly
\[
 \sum_{\substack{a\text{ uncolored}\\\ell_a\geq2}}(\ell_a-2)
 +\sum_{a\text{ colored}}(\ell_a-1)+k_j-1
       =\sum_a(\ell_a-1)+\#\{a:s_a\ne0\}-1
       =c_j-1.                                      \tag{I16.3}
\]
Independence here is on the direct sum of the local jet spaces:
the derivative coefficients at different points and the
$k_j-1$ differences of their constant coefficients are independent.

For each nonreflection value retain jets of length $\ell-1$
at its uncolored ramification points and jets of length $\ell$
at its colored points.  Reflection values have no constraints
and need no retained jets.  This defines an actual effective
divisor $D_{\mathrm{jet}}$ on $Y$, with
\begin{equation}
 \deg D_{\mathrm{jet}}
       =\sum_{c_j\geq2}c_j=N+\nu\leq2N.             \tag{I16.4}
\end{equation}
The direct sum of (I16.3) gives a surjective linear map
$H^0(E|_{D_{\mathrm{jet}}})\to\mathbf C^N$.
Compose it with restriction to obtain the explicit global jet map
\begin{equation}
 J:H^0(Y,E)\longrightarrow H^0(E|_{D_{\mathrm{jet}}})
                         \longrightarrow\mathbf C^N.
                                                       \tag{I16.5}
\end{equation}
Coordinates change its presentation by invertible local operations;
its kernel is exactly the intrinsic first-order equiramification
and common-value condition just proved.

\paragraph{2. Exact forgetful rank.}
Let $F$ be the morphism that takes the original $G$-cover to
its cyclic quotient $X\to Y$, forgetting $h$.  The cyclic root
has no continuous choices over fixed $(Y,R)$, since multiplication
by $d$ on the Picard variety is etale.  The local normal forms
above therefore describe exactly the tangent vectors in a fibre
of $F$.  Composition with the fixed cyclic cover retains the
colored cycles; taking its Galois closure with the existing finite
monodromy marking gives the original Nielsen component locally.
Conversely every deformation of the original cover induces these
local forms in its quotient tower.  This proves the identification
of the tangent fibre, rather than only a numerical comparison.

Infinitesimal postcomposition by $\operatorname{PGL}_2$ is a
three-dimensional subspace of $\ker J$: a target vector field has
equal values over each branch value and no prohibited lower jets
after substitution $h=z^\ell$.  It injects because $h$ is surjective.
There are no infinitesimal automorphisms of $(Y,R)$, since it is
stable in the present signature range.  Since $m\geq n+4$ and
$g\leq n+1$, Serre duality gives $H^1(E)=0$, so
$h^0(E)=2m-g+1$.  We have proved the exact formula
\begin{equation}
 \dim\ker(dF)=2m-g-2-\operatorname{rank}J.          \tag{I16.6}
\end{equation}
The Hurwitz component is locally a finite cover of the configuration
of its $b$ distinct branch values, modulo $\operatorname{PGL}_2$;
its tangent dimension is $b-3$.  Subtracting (I16.6) gives
\begin{equation}
 \boxed{\quad
 \operatorname{rank}(dF)=2n+g-1-N+\operatorname{rank}J.
 \quad}                                             \tag{I16.7}
\end{equation}
This holds at each point, for arbitrary local profiles, with the
actual jet map (I16.5).

If $g\leq n-2$, HD-I12 proves that the cyclic period differential
is injective everywhere, independently of the auxiliary map $h$.
The polarized isogeny of HD-I02 then gives the stronger pointwise
identity
\begin{equation}
 \operatorname{rank}(d\mathrm{Per})
       =2n+g-1-N+\operatorname{rank}J\qquad(g\leq n-2).
                                                       \tag{I16.8}
\end{equation}
For all genera the exact composite remains available: if $\mu$
denotes the primitive multiplication of HD-I05, its remaining
kernel in the cyclic deformation space is $(\operatorname{im}\mu)^\perp$
under Serre duality, and
\[
 \operatorname{rank}(d\mathrm{Per})
 =\operatorname{rank}(dF)
   -\dim\bigl(\operatorname{im}(dF)
                     \cap(\operatorname{im}\mu)^\perp\bigr).
\]

\paragraph{3. The stated nonreflection range.}
Under (I16.2) we have
$\deg(E(-D_{\mathrm{jet}}))\geq2g-1>2g-2$, so
\[
 H^1(E(-D_{\mathrm{jet}}))=0.
\]
The exact restriction sequence therefore makes
$H^0(E)\to H^0(E|_{D_{\mathrm{jet}}})$ surjective.  All $N$ independent
local constraints are globally attainable, and (I16.5) has rank $N$.
Formula (I16.7) proves that $dF$ is onto the full
$(2n+g-1)$-dimensional cyclic deformation space at every point.
Thus each original component has a locally open image in the
cyclic-cover space with its actual colored branching data.

HD-I12 gives full cyclic rank everywhere for $g\leq n-2$;
HD-I13 gives it on a nonempty open subset of each cyclic component
for $g=n-1$; and the fixed-branching connectedness and rank
certificates used in HD-I14 give the same conclusion for $g=n,n+1$.
The open image of every component therefore meets the corresponding
full-rank locus.  The universal factorization bound of HD-I04 is
$2n+g-1$, proving the claimed generic rank under (I16.2).
No generic-position assumption on a pencil or on its ramification
points has been added to the given numerical criterion.

\paragraph{Checks and remaining scope.}
The author, GPT-6 Astra Pro with inherited effort, read the complete
proof on 21 November 2026.  The checks distinguish a movable uncolored
ramification center from a fixed colored center, retain all common
branch-value equations, count their independent local coefficients,
and divide by the actual target automorphisms.  The restriction
sequence proves the smoothness criterion, and the precise cyclic
factor proves the period-rank assertion.  Outside (I16.2), the
explicit finite-dimensional map $J$ still gives (I16.7), and gives
the period rank itself by (I16.8) in the lower-genus range.  Its rank
for every remaining arbitrary Nielsen class is not evaluated here;
the rigid examples HD-I15 have period rank zero and show why this
constraint map cannot be discarded.
